\section{Conclusion}
\label{sec:conclusion}

This paper presented ApproxPilot-SE, a co-evolving surrogate-assisted
DSE framework for approximate accelerators. Search-selected
configurations provide new PPA and QoR measurements that update both
surrogates before the next DSE round. This feedback directs the
available training-label budget toward Pareto-relevant regions
of the legal approximation space.

The experimental results show that this interaction between surrogate
learning and search improves the quality of the final verified design
set. On DCT76, ApproxPilot-SE achieves a verified 3-D PPA HV Gain of
0.085054 at 90\% SSIM retention, compared with 0.023273 for the
stronger baseline, and retains an HV Gain of 0.067996 at 99\%
retention. Under the same 900-label target-system training budget,
co-evolving optimization improves verified HV Gain over an aligned
one-shot workflow by up to 50.84\%. The timing-domain path readout
further reduces Delay MAE by 7.45\% relative to global max pooling
while improving both Spearman correlation and pairwise ranking
accuracy.

The final verified archive also contains high-quality configurations
that simultaneously improve multiple hardware objectives while
remaining close to Exact output quality. A representative DCT76
design retains 99.68\% of the Exact-design SSIM while reducing Area,
Power, and Delay by 19.73\%, 20.50\%, and 12.71\%, respectively.
These results demonstrate the value of coupling surrogate updates
with search-guided system-level evaluation for high-quality
approximate-accelerator design.
